\documentclass[12pt]{article}

\usepackage{enati-template}

\usepackage{graphicx,url}

%\usepackage[brazil]{babel}   
%\usepackage[latin1]{inputenc}  
%\usepackage[T1]{fontenc}
%usepackage[brazilian,english]{babel}
\usepackage[utf8]{inputenc}

\renewcommand{\figurename}{Figura}
\renewcommand{\tablename}{Tabela}
\renewcommand{\refname}{Referências}

     
\sloppy

\title{Instruções para autores do Encontro Anual de Tecnologia da Informação do Oeste Goiano -- ENATI}

\author{Autor1 de Tal, Fulano de Tal\inst{2}, Autor3 do Template\inst{1}, Ciclano de Tal\inst{3} }


\address{Instituto Federal Goiano (IFGoiano)  Campus Iporá\\
  Av. Oeste, 350 -- Parque União - Iporá - GO - Brasil
\nextinstitute
  Department of Computer Science -- University of XXXXX\\
  City, Country.
\nextinstitute
  Departamento de Sistemas e Computação\\
  Universidade XXXXXX (SIGLA) -- Cidade, UF -- Brazil
  \email{\{autor1, autor3\}@ifgoiano.edu.br}
  \email{fulano.de.tal@email.com, ciclano.de.tal@inf.universidade.br}
}

\begin{document} 

\maketitle

\begin{abstract}
This meta-paper describes the style to be used in articles for publication in Annual Meeting of Information Technology of the West Goiano, based on SBC conferences papers template. For papers in English, you should add just an abstract while for the papers in Portuguese, we also ask for an abstract in Portuguese (``resumo''). In both cases, abstracts should not have more than 10 lines and must be in the first page of the paper.
\end{abstract}
     
\begin{resumo} 
  Este meta-artigo descreve o estilo a ser usado na confecção de artigos para publicação nos anais do Encontro Anual de Tecnologia da Informação do Oeste Goiano, com base no modelo de artigos para conferências da SBC. É solicitada a escrita de resumo e abstract apenas para os artigos escritos em português. Artigos em inglês deverão apresentar apenas abstract. Nos dois casos, o autor deve tomar cuidado para que o resumo (e o abstract) não ultrapassem 10 linhas cada, sendo que ambos devem estar na primeira página do artigo.
\end{resumo}


\section{Informações Gerais}

Todos os artigos submetidos para o ENATI, devem ser escritos em Inglês ou Português. A formatação da página deve ser A4, de coluna única e margens superior 3,5cm, inferior 2,5cm, esquerda e direita 2,5cm, sem cabeçalhos ou rodapés.

A fonte principal deve ser Times New Roman, tamanho 12 pontos, com 6 pontos de espaçamento antes de cada parágrafo. Não deve haver numeração nas páginas.

\section{Primeira Página} \label{sec:firstpage}

A primeira página deve mostrar o título do artigo, o nome e endereço dos autores, o \textit{abstract} e o resumo, este último somente para artigos escritos em Português. O título deve ser centralizado, tamanho 16 pontos, negrito e com 12 pontos de espaçamento antes do parágrafo. Os nomes dos autores deve ser centralizado, com tamanho 12 pontos e negrito. Todos os nomes de autores devem ser apresentados na mesma linha, separado por vírgulas e com 12 pontos de espaçamento antes do parágrafo. Os endereços devem ser centralizados, com tamanho 12 pontos e também com 12 pontos de espaçamento antes do parágrafo. Endereços de e-mail devem ser escritos utilizando a fonte Courier New, tamanho 10 pontos, com 6 pontos de espaçamento antes e depois do parágrafo.

O abstract e o resumo (se for o caso) devem ser escritos com fonte Times New Roman, 12 pontos e indentação 0,8cm de ambos os lados. As palavras Abstract e Resumo devem ser escritas em negrito antes do texto.

\section{Seções e Parágrafos}

Os títulos das seções devem estar em negrito, 13pt, alinhados à esquerda. Deve haver espaçamento de 12pt antes e depois de cada título. A numeração das seções é opcional. O primeiro parágrafo de cada seção não deve ser indentado, enquanto que a primeira linha dos parágrafos subsequentes deve ter endentação de 1,27cm.

\subsection{Subseções}

Os títulos das subseções devem estar em negrito, 12pt, alinhados à esquerda.

\section{Figuras e Legendas}\label{sec:figs}


As legendas de figuras e tabelas devem estar sentralizadas caso ocupem apenas uma linha (Figura~\ref{fig:exampleFig1}), do contrário, devem estar justificadas e com endentação de 0,8cm de ambos os lados, como mostrado na Figura~\ref{fig:exampleFig2}. A fonte da legenda deve ser Helvetica, 10 pontos, negrito e com 6 pontos de espaçamento antes e depois do parágrafo.

\begin{figure}[ht]
\centering
\includegraphics[width=.5\textwidth]{fig1.jpg}
\caption{Uma figura típica}
\label{fig:exampleFig1}
\end{figure}

\begin{figure}[ht]
\centering
\includegraphics[width=.3\textwidth]{fig2.jpg}
\caption{Essa figura traz um exemplo de uma legenda que ocupa mais de uma linha, devendo ser formatada de acordo com o informado na Seção~\ref{sec:figs}.}
\label{fig:exampleFig2}
\end{figure}

Em tabelas, deve-se evitar o uso de cores ou sombreamento, assim como bordas grossas, duplas ou desnecessárias. A legenda da tabela deve ser colocada antes da tabela (veja Tabela ~\ref{tab:exTable1}) e sua fonte deve ser Helvetica, 10 pontos, negrito e com 6 pontos de espaçamento antes e depois do parágrafo.

\begin{table}[ht]
\centering
\caption{Variables to be considered on the evaluation of interaction
  techniques}
\label{tab:exTable1}
\includegraphics[width=.7\textwidth]{table.jpg}
\end{table}

\section{Imagens}

Todas as imagens e ilustrações devem ser preferencialmente preto-e-branco ou em tons de cinza, exceto para publicações que estarão disponíveis eletronicamente, como é o caso do ENATI. A resolução da imagem deve ser em torno de 600dpi para imagens preto-e-branco e 150-300dpi para imagens em tons de cinza.

\section{Referências}

Referências bibliográficas devem ser precisas e uniformes. Recomenda-se a citação entre colchetes, ex. \cite{knuth:84}, \cite{boulic:91}, and \cite{smith:99}.

As referências devem ser listadas utilizando fonte com tamanho 12 pontos, com 6 pontos de espaçamento antes de cada referência. A primeira linha de cada referência não deve ser indentada, enquanto que as subsequentes devem ser indentadas em 0,5cm. Além disso, a seção de Referências não deve ser numerada.

\bibliographystyle{sbc}
\bibliography{enati-template}

\end{document}
